\documentclass[11pt]{article}
\usepackage[margin=1in]{geometry}
\usepackage{amsmath,amssymb,amsthm}
\usepackage[T1]{fontenc}
\usepackage{lmodern}
\usepackage{microtype}
\usepackage[colorlinks=true,linkcolor=blue,citecolor=blue,urlcolor=blue]{hyperref}
\hypersetup{pdftitle={A finite degree-elevation bound for central Bernstein polynomials},pdfauthor={GPT-6 Astra}}
\newtheorem{theorem}{Theorem}
\newtheorem{lemma}[theorem]{Lemma}
\newtheorem{corollary}[theorem]{Corollary}
\newcommand{\N}{\mathbb N}
\title{A finite degree-elevation bound for\protect\\central Bernstein polynomials}
\author{GPT-6 Astra}
\date{}
\begin{document}
\maketitle
\begin{abstract}
Let $n$ be a positive even integer and let
$f(x)=k\binom{n}{n/2}x^{n/2}(1-x)^{n/2}$, where $k>1$ and
$h=f(1/2)<1$. We prove that all Bernstein coefficients of $f$ belong to
$[0,1]$ after at most $\lceil nh^2/(1-h^2)\rceil$ degree elevations.
This establishes the finite upper bound asked for by Peter Occil.
The proof uses Sawin's explicit coefficient formula and a telescoping
product inequality. A parity refinement gives
$r_{\min}\le 2\lceil nh^2/(2(1-h^2))\rceil-1$.
\end{abstract}

\section{The question and the result}
For an integer $N\ge0$, the degree-$N$ Bernstein basis is
\[
B_{i,N}(x)=\binom Ni x^i(1-x)^{N-i},\qquad 0\le i\le N.
\]
Every polynomial of degree at most $N$ has a unique expansion in this
basis. Increasing $N$ while keeping the polynomial fixed is called
\emph{degree elevation}.

Fix a positive even integer $n=2m$ and a real number $k>1$, and suppose
\[
h=\frac{k\binom{2m}{m}}{4^m}\in(0,1).
\]
The polynomial under consideration is
\[
f(x)=k\binom{2m}{m}x^m(1-x)^m
     =4^m h x^m(1-x)^m.
\]
Its maximum on $[0,1]$ is $h$, attained at $1/2$. For an integer $r\ge0$,
write
\[
f(x)=\sum_{i=0}^{2m+r}b_i^{(r)}B_{i,2m+r}(x),
\]
and let $r_{\min}$ be the least $r$ for which all these coefficients
belong to $[0,1]$. At $r=0$ the only nonzero coefficient is $k$, so
$r_{\min}\ge1$ whenever it exists.

Occil asked whether $\lfloor nh^2/(1-h^2)\rfloor+1$ elevations always
suffice \cite{Occil}. Sawin derived the elevated coefficients and
located their maximum; a central-binomial asymptotic then gives the
scale $nh^2/(1-h^2)$ \cite{Sawin}. The argument below turns that scale
into a finite bound, uniformly in all the parameters.

\begin{theorem}\label{thm:main}
Under the hypotheses above, put $T=nh^2/(1-h^2)$. Then
\[
1\le r_{\min}\le 2\left\lceil\frac T2\right\rceil-1
                  \le\lceil T\rceil
                  \le\lfloor T\rfloor+1.
\]
At the stated sufficient elevations every coefficient is strictly
less than $1$. Moreover, $r_{\min}$ is odd.
\end{theorem}

\section{The maximum coefficient and its bound}
We give the coefficient calculation in full, including the parity
needed for the finite bound.

\begin{lemma}[Coefficient formula and maximum]\label{lem:coeff}
For $r\ge0$, the nonzero Bernstein coefficients are
\begin{equation}\label{eq:coeff}
b_{m+j}^{(r)}=4^m h\frac{\binom rj}{\binom{2m+r}{m+j}},
\qquad 0\le j\le r.
\end{equation}
Their maximum is attained at $j=r/2$ if $r$ is even, and at
$j=(r-1)/2,(r+1)/2$ if $r$ is odd. For $r\ge1$, set
$R=2\lceil r/2\rceil$. Then
\begin{equation}\label{eq:product}
\max_i b_i^{(r)}
=h\prod_{\ell=1}^{m}\frac{R+2\ell}{R+2\ell-1}.
\end{equation}
\end{lemma}

\begin{proof}
Multiplication by $(x+(1-x))^r=1$ gives
\[
f(x)=4^m h\sum_{j=0}^r\binom rj x^{m+j}(1-x)^{m+r-j},
\]
which proves \eqref{eq:coeff}; all other coefficients are zero.
For $0\le j<r$, the ratio of consecutive nonzero coefficients is
\[
\frac{b_{m+j+1}^{(r)}}{b_{m+j}^{(r)}}
=\frac{(r-j)(m+j+1)}{(j+1)(m+r-j)}.
\]
The numerator minus the denominator on the right is $m(r-2j-1)$.
This establishes the asserted maxima.

If $r=2s$ is even, its maximum divided by $h$ is
\[
4^m\frac{\binom{2s}{s}}{\binom{2s+2m}{s+m}}
=\prod_{\ell=1}^m\frac{2s+2\ell}{2s+2\ell-1},
\]
by cancellation of factorials. If $r=2s-1$ is odd, with $s\ge1$,
use
\[
\binom{2s-1}{s-1}=\frac12\binom{2s}{s},\qquad
\binom{2s+2m-1}{s+m-1}=\frac12\binom{2s+2m}{s+m}.
\]
Thus the maximum at elevation $2s-1$ equals the maximum at elevation
$2s$, proving \eqref{eq:product} in both cases.
\end{proof}

\begin{lemma}[A telescoping bound]\label{lem:bound}
For every integer $r\ge1$ and $R=2\lceil r/2\rceil$,
\begin{equation}\label{eq:bound}
\left(\max_i b_i^{(r)}\right)^2
<h^2\frac{R+n}{R}\le h^2\frac{r+n}{r}.
\end{equation}
\end{lemma}

\begin{proof}
For every real $a>2$,
\[
\left(\frac a{a-1}\right)^2<\frac a{a-2},
\]
since $(a-1)^2=a(a-2)+1$. Here $R\ge2$, so this inequality applies
to every $a=R+2\ell$ in \eqref{eq:product}. Consequently
\[
\prod_{\ell=1}^m\left(\frac{R+2\ell}{R+2\ell-1}\right)^2
<\prod_{\ell=1}^m\frac{R+2\ell}{R+2\ell-2}
=\frac{R+2m}{R}.
\]
Multiplication by $h^2$ proves the strict inequality in
\eqref{eq:bound}. The remaining inequality uses $R\ge r>0$.
\end{proof}

\begin{proof}[Proof of Theorem~\ref{thm:main}]
Let $q=\lceil T/2\rceil$, so $q\ge1$, and take $r=2q-1$. Then
$R=2q\ge T$, whence
\[
h^2\frac{R+n}{R}\le1
\quad\Longleftrightarrow\quad
R\ge\frac{nh^2}{1-h^2}.
\]
Lemma~\ref{lem:bound} now gives $\max_i b_i^{(r)}<1$.
Every coefficient is nonnegative by \eqref{eq:coeff}. This proves
existence of $r_{\min}$ and the first upper bound, including the case
$R=T$ because the product inequality is strict.

To compare the integer bounds, set $L=\lceil T\rceil$. If $L=2a$,
then $2\lceil T/2\rceil-1\le2a-1\le L$. If $L=2a+1$,
then $2\lceil T/2\rceil-1\le2a+1=L$.
Finally $L\le\lfloor T\rfloor+1$. Each of these two larger elevations
satisfies $r\ge T$, so \eqref{eq:bound} gives strict coefficient
bounds at those elevations as well.

The case $r=0$ fails because $k>1$. Every positive even elevation
$2s$ has exactly the same maximum as elevation $2s-1$, by
Lemma~\ref{lem:coeff}. Since all coefficients are nonnegative, the
first successful elevation must therefore be odd.
\end{proof}

\section{An exact criterion and examples}
The product formula also describes the answer exactly, even when the
simple sufficient bound is larger than the minimum.

\begin{corollary}\label{cor:exact}
Let $s_*$ be the least positive integer $s$ satisfying
\[
h\prod_{\ell=1}^m\frac{2s+2\ell}{2s+2\ell-1}\le1.
\]
Then $r_{\min}=2s_*-1$.
\end{corollary}

\begin{proof}
Existence follows from Theorem~\ref{thm:main}. The displayed product
times $h$ is the common maximum at elevations $2s-1$ and $2s$.
All coefficients are nonnegative, and elevation zero fails.
The criterion and the formula for $r_{\min}$ follow.
\end{proof}

For example, take $n=2$ and $h=3/4$. Then $k=3/2$ and $T=18/7$.
At elevation $r=1$, the four coefficients are $(0,1,1,0)$, so
$r_{\min}=1$. The requested bound gives $r_{\min}\le3$.
This example also shows that the asymptotic scale $T$ need not be
a lower bound for the exact minimum.

For an integer-boundary example, take $n=32$ and $h=3/5$. These
parameters satisfy $k>1$: indeed
$\binom{32}{16}/2^{32}=601080390/4294967296<3/5$.
Here $T=18$, and the parity bound gives $r_{\min}\le17$.
The strict inequality in Lemma~\ref{lem:bound} applies with $R=T=18$,
so this boundary case requires no limiting argument.

The proof is entirely algebraic. The accompanying exact-arithmetic
program checks the coefficient formula against repeated degree
elevation, the maximizing indices, the product identity, and the
stated bounds on finite sets of parameters. These checks provide
reproducible verification of the formulas used in the proof.

\begin{thebibliography}{9}
\bibitem{Occil}
Peter Occil,
\emph{On the degree elevation needed to bring Bernstein coefficients
to $[0,1]$}, MathOverflow question 381419,
17 January 2021, revised the same day.
\href{https://mathoverflow.net/q/381419}{mathoverflow.net/q/381419}.

\bibitem{Sawin}
Will Sawin,
Answer to \emph{On the degree elevation needed to bring Bernstein
coefficients to $[0,1]$}, MathOverflow,
17 January 2021, edited 18 January 2021.
\href{https://mathoverflow.net/a/381465}{mathoverflow.net/a/381465}.
\end{thebibliography}
\end{document}
